% Options for packages loaded elsewhere
\PassOptionsToPackage{unicode}{hyperref}
\PassOptionsToPackage{hyphens}{url}
\documentclass[
  ignorenonframetext,
]{beamer}
\newif\ifbibliography
\usepackage{pgfpages}
\setbeamertemplate{caption}[numbered]
\setbeamertemplate{caption label separator}{: }
\setbeamercolor{caption name}{fg=normal text.fg}
\beamertemplatenavigationsymbolsempty
% remove section numbering
\setbeamertemplate{part page}{
  \centering
  \begin{beamercolorbox}[sep=16pt,center]{part title}
    \usebeamerfont{part title}\insertpart\par
  \end{beamercolorbox}
}
\setbeamertemplate{section page}{
  \centering
  \begin{beamercolorbox}[sep=12pt,center]{section title}
    \usebeamerfont{section title}\insertsection\par
  \end{beamercolorbox}
}
\setbeamertemplate{subsection page}{
  \centering
  \begin{beamercolorbox}[sep=8pt,center]{subsection title}
    \usebeamerfont{subsection title}\insertsubsection\par
  \end{beamercolorbox}
}
% Prevent slide breaks in the middle of a paragraph
\widowpenalties 1 10000
\raggedbottom
\AtBeginPart{
  \frame{\partpage}
}
\AtBeginSection{
  \ifbibliography
  \else
    \frame{\sectionpage}
  \fi
}
\AtBeginSubsection{
  \frame{\subsectionpage}
}
\usepackage{iftex}
\ifPDFTeX
  \usepackage[T1]{fontenc}
  \usepackage[utf8]{inputenc}
  \usepackage{textcomp} % provide euro and other symbols
\else % if luatex or xetex
  \usepackage{unicode-math} % this also loads fontspec
  \defaultfontfeatures{Scale=MatchLowercase}
  \defaultfontfeatures[\rmfamily]{Ligatures=TeX,Scale=1}
\fi
\usepackage{lmodern}
\ifPDFTeX\else
  % xetex/luatex font selection
\fi
% Use upquote if available, for straight quotes in verbatim environments
\IfFileExists{upquote.sty}{\usepackage{upquote}}{}
\IfFileExists{microtype.sty}{% use microtype if available
  \usepackage[]{microtype}
  \UseMicrotypeSet[protrusion]{basicmath} % disable protrusion for tt fonts
}{}
\makeatletter
\@ifundefined{KOMAClassName}{% if non-KOMA class
  \IfFileExists{parskip.sty}{%
    \usepackage{parskip}
  }{% else
    \setlength{\parindent}{0pt}
    \setlength{\parskip}{6pt plus 2pt minus 1pt}}
}{% if KOMA class
  \KOMAoptions{parskip=half}}
\makeatother
\usepackage{longtable,booktabs,array}
\usepackage{caption}
\captionsetup[table]{skip=6pt}
\newcounter{none} % for unnumbered tables
\usepackage{calc} % for calculating minipage widths
\usepackage{caption}
% Make caption package work with longtable
\makeatletter
\def\fnum@table{\tablename~\thetable}
\makeatother
\setlength{\emergencystretch}{3em} % prevent overfull lines
\providecommand{\tightlist}{%
  \setlength{\itemsep}{0pt}\setlength{\parskip}{0pt}}
% Manuscript macro/package fallbacks (newcommand -> providecommand).
% Core mathematical packages
\usepackage{amsmath,amssymb,amsthm}
\usepackage{mathtools}
% Permit bibliography and theorem prose to break without overfull boxes.
% Algorithm formatting
% Tables
\usepackage{booktabs}
\usepackage{multirow}
% Graphics
% Cross-referencing with smart naming
% Configure cleveref for custom environments
% Theorem environments with consistent numbering
% Code listings used by supplementary implementation examples
\usepackage{listings}
\lstset{basicstyle=\ttfamily\small,breaklines=true,columns=fullflexible}
% End of theorem environment declarations
% Math operators
\DeclareMathOperator*{\argmax}{arg\,max}
\DeclareMathOperator*{\argmin}{arg\,min}
\DeclareMathOperator{\sign}{sign}
\DeclareMathOperator{\KL}{KL}
\DeclareMathOperator{\Tr}{Tr}
% Custom commands for notation consistency
\providecommand{\calA}{\mathcal{A}}
\providecommand{\calB}{\mathcal{B}}
\providecommand{\calC}{\mathcal{C}}
\providecommand{\calD}{\mathcal{D}}
\providecommand{\calF}{\mathcal{F}}
\providecommand{\calG}{\mathcal{G}}
\providecommand{\calH}{\mathcal{H}}
\providecommand{\calI}{\mathcal{I}}
\providecommand{\calM}{\mathcal{M}}
\providecommand{\calO}{\mathcal{O}}
\providecommand{\calP}{\mathcal{P}}
\providecommand{\calS}{\mathcal{S}}
\providecommand{\calT}{\mathcal{T}}
\providecommand{\calW}{\mathcal{W}}
% Note: \Phi is a standard LaTeX command, do not redefine
% Trust notation
\providecommand{\trust}[2]{\mathcal{T}_{#1 \to #2}}
\providecommand{\trustt}[3]{\mathcal{T}_{#1 \to #2}^{#3}}
% Belief notation
\providecommand{\belief}[2]{\mathcal{B}_{#1}(#2)}
\providecommand{\belieft}[3]{\mathcal{B}_{#1}^{#2}(#3)}
% Cognitive state
\providecommand{\cogstate}[1]{\sigma_{#1}}
% Attack notation
\providecommand{\attack}[1]{\mathcal{A}_{#1}}
\providecommand{\adversary}[1]{\Omega_{#1}}
% Defense notation
\providecommand{\firewall}{\mathcal{F}}
\providecommand{\sandbox}{\mathcal{S}_{box}}
% Probability and expectation
\providecommand{\E}{\mathbb{E}}
\providecommand{\Prob}{\mathbb{P}}
\providecommand{\indicator}{\mathbb{1}}
% QED symbol
\providecommand{\qedsymbol}{$\blacksquare$}
% Page Layout (Slightly smaller margins)
% Hyperlink Styling
% List of Figures and List of Tables in TOC

% Slide layout defaults for warning-clean scientific decks.
\usepackage{etoolbox}
\IfFileExists{xurl.sty}{\usepackage{xurl}}{}
\IfFileExists{seqsplit.sty}{\usepackage{seqsplit}}{\newcommand{\seqsplit}[1]{#1}}
\protected\def\breakseq#1{\seqsplit{#1}}
\protected\def\breaktt#1{\begingroup\ttfamily\seqsplit{#1}\endgroup}
\setlength{\emergencystretch}{6em}
\tolerance=5000
\hbadness=10000
\hfuzz=1pt
\setlength{\tabcolsep}{2pt}
\AtBeginEnvironment{longtable}{\tiny\renewcommand{\arraystretch}{0.86}\setlength{\tabcolsep}{1pt}}
\AtBeginEnvironment{tabular}{\tiny\renewcommand{\arraystretch}{0.86}\setlength{\tabcolsep}{1pt}}
\AtBeginEnvironment{equation}{\tiny}
\AtBeginEnvironment{equation*}{\tiny}
\AtBeginEnvironment{align}{\tiny}
\AtBeginEnvironment{align*}{\tiny}
\AtBeginEnvironment{itemize}{\footnotesize}
\AtBeginEnvironment{enumerate}{\footnotesize}
\AtBeginEnvironment{description}{\footnotesize}
\setbeamerfont{caption}{size=\tiny}
\setbeamerfont{caption name}{size=\tiny}
\setbeamerfont{normal text}{size=\small}
\setbeamerfont{frametitle}{size=\small}
\setbeamerfont{section title}{size=\footnotesize}
\setbeamerfont{subsection title}{size=\footnotesize}
\setbeamertemplate{section page}{%
  \centering
  \begin{beamercolorbox}[sep=12pt,center,wd=\paperwidth]{section title}
    \parbox{0.86\paperwidth}{\centering\usebeamerfont{section title}\insertsection\par}
  \end{beamercolorbox}
}
\setbeamertemplate{subsection page}{%
  \centering
  \begin{beamercolorbox}[sep=8pt,center,wd=\paperwidth]{subsection title}
    \parbox{0.86\paperwidth}{\centering\usebeamerfont{subsection title}\insertsubsection\par}
  \end{beamercolorbox}
}
\setlength{\abovecaptionskip}{2pt}
\setlength{\belowcaptionskip}{0pt}

% Natbib and cross-reference fallbacks — slides don't load natbib
% or cleveref, but combined-PDF manuscript prose may emit these
% commands. The fallback renders citations as a bracketed key list
% and cross-references as detokenized labels so slides don't fail on
% undefined control sequences or raw underscores. \providecommand is
% a no-op if packages load later.
\providecommand{\citep}[1]{[#1]}
\providecommand{\citet}[1]{#1}
\providecommand{\citealp}[1]{#1}
\providecommand{\citeauthor}[1]{#1}
\providecommand{\citeyear}[1]{#1}
\providecommand{\cref}[1]{\texttt{\detokenize{#1}}}
\providecommand{\Cref}[1]{\texttt{\detokenize{#1}}}
\providecommand{\crefrange}[2]{\texttt{\detokenize{#1}}--\texttt{\detokenize{#2}}}
\providecommand{\Crefrange}[2]{\texttt{\detokenize{#1}}--\texttt{\detokenize{#2}}}
\providecommand{\paragraph}[1]{\textbf{#1}\ }

% Additional theorem-like environments from manuscript preamble.
\newtheorem{observation}{Observation}
\newtheorem{formalization}{Formalization}
\newtheorem{principle}{Principle}

% Environment-providing packages carried over from the manuscript.
\usepackage{algpseudocode}

% Non-floating stand-in for the `algorithm` float.
\newenvironment{algorithm}[1][]{%
  \par\medskip\noindent\rule{\linewidth}{0.4pt}\par\nobreak\small
  \renewcommand{\caption}[1]{\par\noindent\textbf{##1}\par}}{%
  \par\nobreak\noindent\rule{\linewidth}{0.4pt}\par\medskip}

\newtheorem{proposition}{Proposition}
\newtheorem{hypothesis}{Hypothesis}
\newtheorem{remark}[theorem]{Remark}
\newtheorem{axiom}{Axiom}
\newtheorem{property}{Property}
\usepackage{bookmark}
\IfFileExists{xurl.sty}{\usepackage{xurl}}{} % add URL line breaks if available
\urlstyle{same}
% fallback for those not using the hyperref driver hyperxmp:
\makeatletter
\@ifundefined{xmpquote}{\newcommand{\xmpquote}[1]{#1}}{}
\makeatother
\hypersetup{
  hidelinks,
  pdfcreator={LaTeX via pandoc}}

\author{\texorpdfstring{}{}}
\date{}

\begin{document}

\begin{frame}
\newpage
\end{frame}

\section{Discussion: Theoretical Implications, Limitations, and Future
Directions}\label{sec:discussion}

\begin{frame}[allowframebreaks]{Discussion: Theoretical Implications,
Limitations, and Future Directions}
This section examines the theoretical implications of the Cognitive
Integrity Framework (\cref{sec:theoretical-implications}), formal
limitations (\cref{sec:formal-limitations}) and boundary conditions
(\cref{sec:limitations}), relationship to prior work
(\cref{sec:related-work}), governance implications
(\cref{sec:governance}), and future research directions
(\cref{sec:future-directions}).
\end{frame}

\subsection{Empirical Validation Summary (Part
2)}\label{empirical-validation-summary-part-2}

\begin{frame}[allowframebreaks]{Empirical Validation Summary (Part 2)}
The formal mechanisms proposed in this paper have been validated through
extensive architecture-aware parametric simulation in \textbf{Part 2:
Computational Validation} \cite{friedman2026cogsec2}. In experiments
across the four production multiagent architectures Part 2 evaluates
(Claude Code, AutoGPT, CrewAI, LangGraph), the parametric simulation
establishes a design-level detection ceiling of \textbf{96--100\%}
against a corpus of 950 cognitive attacks; the prototype pipeline
achieves a mean detection rate of \textbf{86.3\%} {[}95\% CI: 85.5\%,
87.1\%{]} across 30 random seeds (Part 2, Results).

\end{frame}

\begin{frame}[allowframebreaks]{Empirical Validation Summary (Part 2)}
Key empirical findings include:

\end{frame}

\begin{frame}[allowframebreaks]{Empirical Validation Summary (Part 2)}
\begin{enumerate}
\tightlist
\item
  \textbf{Defense Composition}: Consistent with the series and parallel
  detection-rate theorems (\cref{thm:series-detection},
  \cref{thm:parallel-detection}), layered defenses composed
  multiplicatively, with the full framework outperforming any single
  mechanism. Part 2's ablation studies qualify this: on their corpus the
  composition is dominated rather than balanced, with the Behavioral
  Invariants module accounting for almost all of the composed pipeline's
  detection and every other module measuring at or near zero
  \emph{marginal} contribution --- what each adds beyond the rest of the
  stack, which is not the same quantity as what each detects alone.
\item
  \textbf{Trust Boundedness}: The \(\delta^d\) trust decay parameter
  successfully prevented ``trust laundering'\,' amplification attacks in
  all tested topologies, validating Formula 4.
\item
  \textbf{Architecture Dependence}: As predicted by our threat model,
  the architecture with the lowest parametric detection in Part 2 is
  AutoGPT's autonomous-mesh topology, whose distributed trust structure
  offers the widest lateral-movement surface; hierarchical topologies
  concentrate risk in orchestrator compromise instead. No peer-to-peer
  architecture was evaluated.
\item
  \textbf{Adversarial Training}: In Part 2's closed-form
  adversarial-training design model, five rounds on the Claude Code
  architecture raise hardened detection from 52.0\% (Round 1) to 67.9\%
  (Round 5), \(+23.2\) pp over the pre-AT baseline. The per-round gain
  sequence \((7.3, 5.6, 5.0, 2.9, 2.4)\) pp is approximately linear
  rather than geometrically decaying, and the projected Nash equilibrium
  of 100.0\% is a property of the model's assumed gains, not a measured
  equilibrium. These are design-model values, not pipeline-in-the-loop
  measurements.
\end{enumerate}

\end{frame}

\begin{frame}[allowframebreaks]{Empirical Validation Summary (Part 2)}
These results confirm that the abstract properties proven
here---boundedness, composition, and detectability---translate into
concrete security improvements when implemented in realistic agent
architectures.
\end{frame}

\subsection{Theoretical
Implications}\label{sec:theoretical-implications}

\begin{frame}[allowframebreaks]{Why Composable Defenses Are Necessary}
\protect\phantomsection\label{why-composable-defenses-are-necessary}
The defense composition algebra (\cref{thm:series-detection},
\cref{thm:parallel-detection}) formalizes a principle implicit in
security practice: layered defenses provide multiplicative rather than
additive protection. Each defense mechanism addresses a distinct attack
surface:

\end{frame}

\begin{frame}[allowframebreaks]{Why Composable Defenses Are Necessary}
\begin{table}[htbp]
\centering
\caption{Defense mechanisms and their target attack surfaces.}
\label{tab:defense-surfaces}
\begin{tabular}{@{}ll@{}}
\toprule
Defense Layer & Target Attack Surface \\
\midrule
Cognitive Firewall & Input-based injection \\
Belief Sandbox & Unverified content propagation \\
Tripwires & Belief manipulation \\
Trust Calculus & Delegation abuse \\
Byzantine Consensus & Coordination attacks \\
\bottomrule
\end{tabular}
\end{table}

\end{frame}

\begin{frame}[allowframebreaks]{Why Composable Defenses Are Necessary}
The orthogonality of these surfaces explains why no single mechanism
suffices: an attack that bypasses input filtering may still violate
behavioral invariants; an attack that evades pattern matching may still
trigger belief drift detection.

Empirical ablation studies in Part 2 test this prediction and return a
sharper answer than a ranked list. On the 100-attack ablation corpus a
single module dominates: removing Behavioral Invariants costs
\(\Delta\text{TPR} \approx -0.847\) against a full pipeline of
\(\approx 0.959\), the Firewall and Tripwire cost \(\approx -0.010\)
each, and Consensus, Detection, Provenance, Sandbox and Trust Calculus
register no change at all. Those zeros are \emph{marginal} quantities
and are not evidence of zero capability: a module whose detections are a
subset of another's is invisible to leave-one-out removal however much
it catches on its own. What the ablation establishes is therefore weaker
\end{frame}

\begin{frame}[allowframebreaks]{Why Composable Defenses Are Necessary}
than the theorem and more useful than a ranking --- no single mechanism
outperforms the composed ensemble, consistent with
\cref{thm:series-detection}, but on this corpus the composition is
dominated by one factor rather than balanced across many, so the
multiplicative bound holds without being informative.
\end{frame}

\begin{frame}[allowframebreaks]{The Trust Boundedness Guarantee}
\protect\phantomsection\label{the-trust-boundedness-guarantee}
The bounded trust theorem (\cref{thm:trust-bounded}) represents a
structural guarantee against trust amplification attacks. Unlike
detection-based defenses that may be evaded by novel attacks, the
\(\delta^d\) decay bound is algebraic: it holds for any attack type, any
adversary capability, and any delegation chain length. This makes it a
\emph{formal} rather than \emph{empirical} security property.

\end{frame}

\begin{frame}[allowframebreaks]{The Trust Boundedness Guarantee}
The decay factor \(\delta \in [0, 1)\) creates a tradeoff:

\end{frame}

\begin{frame}[allowframebreaks]{The Trust Boundedness Guarantee}
\begin{itemize}
\tightlist
\item
  Lower \(\delta\): Stronger security, limited delegation utility
\item
  Higher \(\delta\): More delegation flexibility, weaker bounds
\end{itemize}

\end{frame}

\begin{frame}[allowframebreaks]{The Trust Boundedness Guarantee}
Organizations must calibrate this tradeoff based on their threat model
(\cref{sec:operator-posture} in Part 3). Domain-calibrated \(\delta\)
choices --- from millisecond OODA cycles in drone swarms to year-scale
diplomatic agents --- are analyzed in Part 3 \cite{friedman2026cogsec3}
across ten critical operational domains.
\end{frame}

\begin{frame}[allowframebreaks]{Information-Theoretic Detection Limits}
\protect\phantomsection\label{information-theoretic-detection-limits}
The stealth-impact fundamental limit (\cref{thm:stealth-impact})
establishes that certain attacks are \emph{provably} undetectable
without unacceptable false positive rates. This is not a limitation of
our specific mechanisms but a fundamental bound analogous to Shannon's
channel capacity---some attacks simply cannot be detected without
additional information.

This has practical implications: security architectures should not
promise detection of all possible attacks. Instead, they should
characterize the detection boundary and provide containment for attacks
that cross it.
\end{frame}

\begin{frame}[allowframebreaks]{Architecture-Specific Vulnerability
Patterns}
\protect\phantomsection\label{architecture-specific-vulnerability-patterns}
The formal framework reveals why different multiagent architectures
exhibit different vulnerability profiles:

\end{frame}

\begin{frame}[allowframebreaks]{Architecture-Specific Vulnerability
Patterns}
\begin{table}[htbp]
\centering
\caption{Theoretical vulnerability analysis by architecture type.}
\label{tab:architecture-theory}
\begin{tabular}{@{}lll@{}}
\toprule
Architecture & Theoretical Vulnerability & Formal Mitigation \\
\midrule
Hierarchical & Single point of trust concentration & Byzantine-tolerant orchestrator \\
Peer-to-peer & Uniform trust enables lateral movement & Trust decay on delegation \\
Role-based & Logical (not cryptographic) boundaries & Attestation per role transition \\
State machine & State integrity assumption & State hash verification \\
\bottomrule
\end{tabular}
\end{table}

\end{frame}

\begin{frame}[allowframebreaks]{Architecture-Specific Vulnerability
Patterns}
These are structural properties of the architectures themselves, not
implementation-specific weaknesses.
\end{frame}

\subsection{Formal Limitations}\label{sec:formal-limitations}

\begin{frame}[allowframebreaks]{Assumption Dependencies}
\protect\phantomsection\label{assumption-dependencies}
The formal guarantees of CIF depend on specific assumptions. Violation
of these assumptions degrades or eliminates security properties:

\end{frame}

\begin{frame}[allowframebreaks]{Assumption Dependencies}
\begin{table}[htbp]
\centering
\caption{Impact of assumption violations on formal guarantees.}
\label{tab:assumption-violations}
\begin{tabular}{@{}lp{5cm}@{}}
\toprule
Assumption & Guarantee Impacted \\
\midrule
Honest orchestrator & $\Omega_5$ attacks succeed; systemic compromise possible \\
$n \geq 3f+1$ agents & Byzantine consensus fails (\cref{thm:byzantine-req}) \\
Authenticated channels & $\Omega_4$ coordination attacks expand \\
Known attack patterns & Zero-day evasion possible \\
\bottomrule
\end{tabular}
\end{table}
\end{frame}

\begin{frame}[allowframebreaks]{Scalability Constraints}
\protect\phantomsection\label{scalability-constraints}
The formal framework imposes scaling limitations:

\begin{align}
\label{eq:memory-scaling}
M_{\text{trust}} &= O(n^2) \\
\label{eq:consensus-scaling}
L_{\text{consensus}} &= O(n^2)
\end{align}

The quadratic trust matrix (\cref{eq:memory-scaling}) limits practical
deployment to systems with moderate agent counts. Sparse trust
representations or hierarchical trust structures may enable scaling to
larger systems.

Consensus latency (\cref{eq:consensus-scaling}) suggests that Byzantine
consensus should be reserved for critical decisions rather than applied
universally.
\end{frame}

\begin{frame}[allowframebreaks]{Inherent Detection Gaps}
\protect\phantomsection\label{inherent-detection-gaps}
Certain attack types are formally difficult to detect:

\end{frame}

\begin{frame}[allowframebreaks]{Inherent Detection Gaps}
\begin{itemize}
\tightlist
\item
  \textbf{Semantic equivalence}: Attacks preserving meaning while
  changing syntax evade pattern-based detection
\item
  \textbf{Progressive drift}: Sub-threshold changes that accumulate over
  time (\cref{eq:progressive-drift})
\item
  \textbf{Orchestrator compromise}: Outside the
  \(\Omega_1\)--\(\Omega_4\) threat model
\end{itemize}

\end{frame}

\begin{frame}[allowframebreaks]{Inherent Detection Gaps}
These are not implementation failures but formal limitations of the
detection paradigm.
\end{frame}

\subsection{Relationship to Prior Work}\label{sec:related-work}

\begin{frame}[allowframebreaks]{Relationship to Prior Work}
CIF builds on and extends several research traditions:

\textbf{Byzantine Fault Tolerance}: Classical BFT (PBFT, etc.) assumes
crash or arbitrary faults. CIF extends this to cognitive
manipulation---agents that appear functional but hold corrupted beliefs.

\textbf{Trust Management Systems}: Prior systems (PolicyMaker, SPKI,
etc.) focus on authorization decisions. CIF addresses continuous trust
evolution with provable decay bounds.

\textbf{AI Safety and Alignment}: Constitutional AI
\cite{bai2022constitutional} and similar approaches address single-agent
alignment. CIF extends these concepts to multi-agent coordination
integrity.

\end{frame}

\begin{frame}[allowframebreaks]{Relationship to Prior Work}
\textbf{Prompt Injection Defenses}: Existing defenses focus on
single-agent scenarios. CIF addresses the propagation and amplification
of attacks across agent networks.

\textbf{Cognitive Science and Active Inference}: The Active Inference
framework \cite{david2021aic} provides a complementary perspective from
cognitive science, modeling agents as entities that minimize prediction
error through continuous perception-action loops. The Active Inference
Conflict (AIC) model extends classical decision frameworks like OODA
loops (Observe-Orient-Decide-Act) by situating conflict as a multiscale
process of communication, trust, and relationship management---themes
that directly inform CIF's trust calculus. AIC's treatment of BOLTS
components (Business, Operations, Legal, Technical, Social) also informs
our analysis of cyberphysical cognitive systems where cognitive security
spans multiple operational domains. Critically, the Active Inference
\end{frame}

\begin{frame}[allowframebreaks]{Relationship to Prior Work}
perspective illuminates why belief manipulation attacks are particularly
dangerous: agents minimizing variational free energy will actively seek
information confirming their current beliefs, creating self-reinforcing
loops when those beliefs are corrupted. CIF's tripwire mechanism
interrupts this by detecting when prediction error patterns deviate from
baseline, analogous to detecting abnormal precision weighting in
cognitive systems.

\textbf{Pattern Languages for Cognitive Security}: The COGSEC ATLAS
\cite{cogsecatlas2023} provides a practitioner-oriented complement to
CIF's formal approach, cataloging 995 cognitive security patterns
organized by type (Vulnerability, Exploit, Remedy, Practice,
Accelerator, Moderator, Condition). Where CIF provides provable
guarantees and formal composition rules, the Atlas offers an
empirically-grounded taxonomy of observed attack patterns and defensive
\end{frame}

\begin{frame}[allowframebreaks]{Relationship to Prior Work}
practices---such as the Devil's Advocate and Key Assumptions Check
techniques for countering groupthink and confirmation bias. The
hierarchical parent-child structure of Atlas patterns maps naturally
onto CIF's adversary class hierarchy, suggesting opportunities for
formal verification of pattern-based defenses using the mechanisms
described in \cref{sec:defense-mechanisms}.

The novel contribution is the integration of these concerns into a
unified formal framework with composable guarantees.
\end{frame}

\subsection{Governance and Policy Implications}\label{sec:governance}

\begin{frame}[allowframebreaks]{The Regulatory Gap}
\protect\phantomsection\label{the-regulatory-gap}
Current AI regulation lacks cognitive security provisions:

\end{frame}

\begin{frame}[allowframebreaks]{The Regulatory Gap}
\begin{table}[htbp]
\centering
\caption{Regulatory gaps for cognitive security.}
\label{tab:regulatory-gaps}
\begin{tabular}{@{}lp{4cm}p{4cm}@{}}
\toprule
Regulation & Current Focus & Cognitive Security Gap \\
\midrule
EU AI Act & Risk classification, transparency & No inter-agent trust requirements \\
NIST AI RMF & Risk management lifecycle & Limited multiagent-specific guidance \\
ISO 42001 & AI management systems & Process-focused, not cognitive-state focused \\
\bottomrule
\end{tabular}
\end{table}
\end{frame}

\begin{frame}[allowframebreaks]{Recommendations for Policy}
\protect\phantomsection\label{recommendations-for-policy}
We propose that regulators consider:

\end{frame}

\begin{frame}[allowframebreaks]{Recommendations for Policy}
\begin{enumerate}
\item \textbf{Cognitive Security Audits}: Mandatory assessment of inter-agent trust mechanisms for high-risk deployments
\item \textbf{Transparency Requirements}: Disclosure of trust hierarchies and delegation policies
\item \textbf{Incident Reporting}: New category for cognitive attacks
\item \textbf{Certification Pathways}: Industry certification for cognitive security practices
\end{enumerate}
\end{frame}

\subsection{Future Theoretical Directions}\label{sec:future-directions}

\begin{frame}[allowframebreaks]{Adaptive Defense Theory}
\protect\phantomsection\label{adaptive-defense-theory}
The detection degradation problem suggests a need for adaptive defenses.
Formal treatment as a game-theoretic equilibrium:

\begin{equation}
\label{eq:adaptive-defense}
\pi^{*}_{\text{defense}} = \argmax_{\pi} \mathbb{E}\left[\sum_t \gamma^t r(s_t, a_t)\right]
\end{equation}

requires solving the partial observability problem---defenders cannot
directly observe attacker intent.
\end{frame}

\begin{frame}[allowframebreaks]{Cross-System Trust Federation}
\protect\phantomsection\label{cross-system-trust-federation}
Extending trust calculus across organizational boundaries:

\begin{equation}
\label{eq:cross-system-trust}
\mathcal{T}_{i \to j}^{\text{cross}} = f(\mathcal{T}_{\text{local}}, \mathcal{T}_{\text{reputation}}, \mathcal{T}_{\text{attestation}})
\end{equation}

The primary challenge is trust calibration---mapping heterogeneous trust
semantics across systems with different threat models.
\end{frame}

\begin{frame}[allowframebreaks]{Emergent Behavior Security}
\protect\phantomsection\label{emergent-behavior-security}
As multiagent systems scale, emergent collective behaviors become
security-relevant. Open questions include:

\end{frame}

\begin{frame}[allowframebreaks]{Emergent Behavior Security}
\begin{itemize}
\item Formal characterization of beneficial vs. malicious emergence
\item Detection of emergent coordination patterns indicating compromise
\item Sandboxing that preserves beneficial emergence while constraining malicious emergence
\end{itemize}

\end{frame}

\begin{frame}[allowframebreaks]{Emergent Behavior Security}
The colony cognitive security perspective developed in
\cref{sec:eusocial-cogsec} provides initial formal foundations for these
questions.
\end{frame}

\begin{frame}[allowframebreaks]{Long-Horizon Agent Security}
\protect\phantomsection\label{long-horizon-agent-security}
Agents operating over extended time horizons (days, weeks, months) face
additional challenges:

\end{frame}

\begin{frame}[allowframebreaks]{Long-Horizon Agent Security}
\begin{itemize}
\item \textbf{Memory integrity}: Verification of accumulated knowledge
\item \textbf{Goal stability}: Distinguishing legitimate evolution from adversarial drift
\item \textbf{Temporal consistency}: Decisions consistent with historical context
\end{itemize}

\end{frame}

\begin{frame}[allowframebreaks]{Long-Horizon Agent Security}
The trust calculus extends naturally to temporal trust:
\(\mathcal{T}^{t} = \mathcal{T}^{t-1} \cdot \delta_{\text{time}}\) where
\(\delta_{\text{time}}\) encodes trust decay over time.
\end{frame}

\begin{frame}[allowframebreaks]{The Cognitive Security Research Agenda}
\protect\phantomsection\label{the-cognitive-security-research-agenda}
We propose a research agenda organized by time horizon:

\textbf{Near-term (1--2 years)}:

\end{frame}

\begin{frame}[allowframebreaks]{The Cognitive Security Research Agenda}
\begin{itemize}
\item Standardized formal verification benchmarks
\item Integration of CIF mechanisms into production frameworks
\item Theoretical analysis of real-world attack patterns
\end{itemize}

\end{frame}

\begin{frame}[allowframebreaks]{The Cognitive Security Research Agenda}
\textbf{Medium-term (2--5 years)}:

\end{frame}

\begin{frame}[allowframebreaks]{The Cognitive Security Research Agenda}
\begin{itemize}
\item Game-theoretic foundations for adaptive defenses
\item Cross-organizational trust federation protocols
\item Hardware-backed cognitive security guarantees
\end{itemize}

\end{frame}

\begin{frame}[allowframebreaks]{The Cognitive Security Research Agenda}
\textbf{Long-term (5+ years)}:

\end{frame}

\begin{frame}[allowframebreaks]{The Cognitive Security Research Agenda}
\begin{itemize}
\item Formal verification of emergent agent behavior
\item Self-healing cognitive security systems
\item Integration with broader AI safety theory
\end{itemize}

\end{frame}

\begin{frame}[allowframebreaks]{The Cognitive Security Research Agenda}
The formal foundations established in this work---bounded trust,
composable defenses, information-theoretic limits---provide a stable
basis for this evolving research program.
\end{frame}

\end{document}
